\section{Discussion}
\label{discussion}

\subsection{Interpretation}

The experiments support two conclusions. First, parameter optimization can
materially improve PureCLIP call sets, but the magnitude of improvement depends
on the dataset. Relative to the shared default, nine of 12 LLM runs and 11 of 12
TPE runs found a higher score. Gains ranged from negligible
changes to increases above 0.20. This variation supports dataset-specific tuning,
although the current sample is too small to identify which experimental
properties predict a gain.

Second, biological context influenced LLM performance, but not uniformly. The
positive HNRNPK ablation result is biologically plausible: C-rich motif and
footprint information directed the agent toward sharper, more strongly supported
calls. However, the same prior formulation reduced the best score on PUM2,
U2AF2, and slightly on SRSF1, while SF3B4 was unchanged. For U2AF2, sequence
motifs are an incomplete description of binding because position relative to
splice sites is central. More generally, one trajectory per condition cannot
distinguish a protein-class effect from dataset idiosyncrasies or stochastic LLM
behavior. The HNRNPK result should therefore be treated as a candidate mechanism
for further testing, not evidence that priors generally benefit proteins with
degenerate motifs.

TPE had the higher observed score in ten of 12 completed pairs. This is a
consistent descriptive pattern across both cell lines. It is not, however,
evidence of intrinsic optimizer superiority:
the TPE trajectories usually contain roughly twice as many evaluations. A fair
claim concerns the implemented protocol---TPE more often returned the higher
best score under these stopping rules---rather than matched-budget efficiency.

The study also shows why interpretability must be assessed carefully. The LLM
produced readable rationales and adapted its proposals after unfavorable score
changes. Nevertheless, a plausible explanation is not a causal analysis. Several
parameters changed together, and the objective provides no gradient that assigns
an improvement to one change. Factorial follow-up experiments around selected
LLM proposals would be required to test those explanations.

The exploratory prompt comparison identifies a further methodological
consideration. Although the model received the decomposed objective, some
parameter proposals were not directionally consistent with those components.
Providing more extensive instructions did not eliminate this behavior in the
small set of inspected traces. Moreover, target-gene information could not be
linked to the evaluated call sets because peak-to-gene annotations were absent
from the feedback. These observations are insufficient to compare prompt
variants, but they indicate that an LLM optimizer should be evaluated not only
by its final score but also by the consistency of its proposals with the
information available at each iteration.

\subsection{Limitations of the comparison}

The common pipeline and objective remove an important source of confounding, but
the computational budgets were not matched. LLM runs began at the default and
typically stopped after five to nine evaluations. TPE began with sampled
configurations, retained a ten-trial startup phase, and completed 14--16
evaluations. The observed best scores therefore compare available trajectories
rather than equal search effort. Runtime, API cost, and the number of successful
evaluations should all be fixed or reported in a definitive comparison.

The biological evaluation has further limitations. All experiments were
restricted to chromosome~21 and were run once, so neither genome-wide transfer
nor run-to-run variability was measured. The analyzed subset represents eight of
15 RBPs in the registry and was determined by available completed runs rather
than random or stratified sampling. The registry demonstrates pipeline scope,
not benchmark-wide performance. The reproducibility background uses only
three shuffled interval sets. Mononucleotide shuffling controls base composition
but not dinucleotide structure. ENCODE reference regions are outputs of another
analysis pipeline rather than experimentally verified ground truth. Finally, the
motif component selects the most enriched PWM from a protein-specific catalog.
Catalog size and motif quality may therefore affect comparability among RBPs, and
a sequence-based term is poorly matched to strongly positional binders.

The composite objective also encodes normative choices. Its fixed weights treat
replicates as equally reliable and value motif support identically across
proteins. A high composite may conceal a loss in one component, as observed for
HNRNPK reference recall. For this reason, decomposed metrics should accompany the
aggregate score, and future objectives should incorporate replicate quality and
protein-appropriate evidence.

\subsection{Future work}

The immediate priority is a replicated, equal-budget comparison that extends to
more of the dataset registry. Each optimizer should receive the same number of successful
pipeline evaluations, with multiple TPE seeds and multiple LLM sampling seeds.
Reporting best-score trajectories against both evaluation count and wall-clock
time would separate sample efficiency from computational cost. The prior ablation
should include several RBPs in each pre-specified motif class.

Methodologically, positional binders require a feature based on distance to
annotated splice sites or branch points rather than a generic PWM score.
Dinucleotide-preserving sequence backgrounds and a more robust reproducibility
estimator would strengthen the remaining components. A controlled prompt study
could replay identical optimization histories under pre-specified prompt
variants and quantify whether proposed parameter changes are consistent with the
reported component scores. Peak-to-gene annotations would be required before
target-gene priors could be evaluated. Finally, the best chromosome-21
configurations should be evaluated genome-wide on held-out data;
without this step, the optimized settings should not be interpreted as validated
production parameters.
